%%%%%%%%%%%%%%%%%%%%%%%%%%%%%%%%%%%%%%%%%%%%%%%%%%%%%%%%%%%%%%%%%%%%%%%%%%%%%%%%
%2345678901234567890123456789012345678901234567890123456789012345678901234567890
%        1         2         3         4         5         6         7         8

\documentclass[letterpaper, 10 pt, conference]{ieeeconf}  % Comment this line out if you need a4paper

%\documentclass[a4paper, 10pt, conference]{ieeeconf}      % Use this line for a4 paper

\IEEEoverridecommandlockouts                              % This command is only needed if 
                                                          % you want to use the \thanks command

\overrideIEEEmargins                                      % Needed to meet printer requirements.


% The following packages can be found on http:\\www.ctan.org
%\usepackage{graphics} % for pdf, bitmapped graphics files
%\usepackage{epsfig} % for postscript graphics files
%\usepackage{mathptmx} % assumes new font selection scheme installed
%\usepackage{times} % assumes new font selection scheme installed
%\usepackage{amsmath} % assumes amsmath package installed
%\usepackage{amssymb}  % assumes amsmath package installed

\usepackage{amsmath,amsfonts}
\usepackage{algorithmic}
\usepackage{algorithm}
\usepackage{array}
\usepackage[caption=false,font=normalsize,labelfont=sf,textfont=sf]{subfig}
\usepackage{textcomp}
\usepackage{stfloats}
\usepackage{url}
\usepackage{verbatim}
\usepackage{graphicx}
\usepackage{booktabs} % better hlines for tables
\usepackage{makecell}
\usepackage{cite}
\usepackage[dvipsnames]{xcolor}
\usepackage{lipsum}


% https://tex.stackexchange.com/questions/247104/hyperref-doesnt-link-cite-command hero
\makeatletter
\let\NAT@parse\undefined
\makeatother
\usepackage{multirow}
\usepackage{hyperref}
\usepackage{tabularx}
\usepackage{listings}
\usepackage{placeins}

% Define 'Y' as a centered, stretchable 'X' column
\newcolumntype{Y}{>{\centering\arraybackslash}X}



% Override table caption style to match figure captions (remove small caps)
\newlength{\figurecaptionskip}
\newlength{\tablecaptionskip}
\setlength{\figurecaptionskip}{0pt}
\setlength{\tablecaptionskip}{4pt}

\makeatletter
\long\def\@makecaption#1#2{%
\ifx\@captype\@IEEEtablestring%
\vskip\tablecaptionskip%
{\footnotesize #1.~~ #2\par}%
\vskip\tablecaptionskip%
\else
\vskip\figurecaptionskip%
\setbox\@tempboxa\hbox{\footnotesize #1.~~ #2}%
\ifdim \wd\@tempboxa >\hsize%
\setbox\@tempboxa\hbox{\footnotesize #1.~~ }%
\parbox[t]{\hsize}{\footnotesize \noindent\unhbox\@tempboxa#2}%
\else%
\ifcenterfigcaptions \hbox to\hsize{\footnotesize\hfil\box\@tempboxa\hfil}%
\else \hbox to\hsize{\footnotesize\box\@tempboxa\hfil}%
\fi\fi\fi}
\makeatother

\hypersetup{
    colorlinks,
    linkcolor={red!50!black},
    citecolor={red!50!black},
    urlcolor={blue!80!black}
}


%%%% ANONYMOUS SUBMISSION TOGGLE
% Set to `true` for anonymous submission, `false` to show authors and funding.
\newif\ifanonymous
\anonymousfalse  % <--- change to \anonymoustrue for anonymous submission

%%%% APPENDIX TOGGLE
% Set to `true` to include appendices, `false` to hide them (e.g. for page-limited submissions).
\newif\ifappendixincluded
\appendixincludedtrue  % <--- change to \appendixincludedfalse to hide appendices

% Use \appendixnote{sentence} in the body text to include a sentence only when appendices are present.
\newcommand{\appendixnote}[1]{\ifappendixincluded#1\fi}


%%%% COMMANDS

\newcommand{\thomas}[1]{{\color{ForestGreen} TL: #1}}
\newcommand{\thomasRemark}[1]{{\color{ForestGreen} {\bf TL: #1}}}

\newcommand{\marco}[1]{{\color{Purple} MM: #1}}

\title{\LARGE \bf
On the Generalization Capabilities, Design Choices and Limitations of Keypoint Imitation Learning}

\ifanonymous
\author{Anonymous$^{*}$ - Double Blind Submission
\thanks{*Affiliations anonymized for double-blind submission}%
}
\else
\author{Thomas Lips$^{*,1}$, Marco Moletta$^{*,2}$, Michael C. Welle$^{2,3}$, Danica Kragic$^{2}$ and Francis wyffels$^{1}$% <-this % stops a space
\thanks{*Shared first authors. thomas.lips@UGent.be, moletta@kth.se}% <-this % stops a space
\thanks{$^{1}$ AI and Robotics Lab, IDLAB-AIRO, Ghent University-imec}%
\thanks{$^{2}$ Robotics, Perception and Learning Lab, (RPL), EECS, KTH Royal Institute of Technology.}%
\thanks{$^{3}$ INCAR Robotics AB, Stockholm, Sweden. }%
}
\fi


\begin{document}



\maketitle
\thispagestyle{empty}
\pagestyle{empty}


\begin{abstract}
%context
RGB-based imitation learning requires many demonstrations to generalize to unseen objects or scenes, motivating research into intermediate representations to improve generalization for robotic manipulation.
%problem
Visual foundation models enable one-shot extraction of keypoints to provide such representation. However, it remains unclear how to integrate them into imitation learning optimally and when they outperform alternative representations.
%what and results
We combine approaches from previous works on keypoint imitation learning (KIL) and investigate several design choices to provide practical guidelines. Using over 2000 real-world rollouts, we also assess the generalization capabilities of KIL to unseen objects and scene variations. KIL achieves a 75\% overall success rate across five tasks, significantly outperforming the RGB baseline (47\%) and performing on par with S$^2$-diffusion (73\%). Finally, we explore the limitations of the foundation models used for keypoint extraction and extend KIL to tasks with multiple object instances.
%outlook
Our results confirm KIL as a data-efficient approach for robot learning, though it does not outperform alternative representations and inherits limitations of the foundation models used for keypoint extraction. 
\end{abstract}


\section{Introduction}
\label{sec:introduction}


% We want robots that can operate in human enviroments. such environments have both large diversity and change continuously. This requires robots to handle a large diversity of objects. Obtaining such generalization is challenging and often requires data-driven methods combined with large amounts of data. Keypoints can help to make this more efficient by conditioning action generation on a low-dimensional set of task-relevant and semantically consistent points (cite papers).

% The contributions of this paper are as follows:
% \begin{itemize}
%     \item 
% \end{itemize}

%context
Robotic manipulation aims to enable robots to reliably execute physical tasks in unstructured environments, requiring generalization to different poses, unseen objects and scene variations. Imitation learning is a promising approach to achieve this~\cite{chi2025diffusion_policy,barreiros2025LBM-1}, but RGB-based policies still require large amounts of data to handle the aforementioned task variations\cite{barreiros2025LBM-1,lin2025imitation_scaling_laws}. This motivates research into intermediate representations for improved generalization~\cite{ze20243d,yang2025S2,zhu2023viola,zhu2023learning,wang2025skil,di2024KAT,Levy2025P3PO}.

%problem
Keypoints are an attractive option for such representation: they are compact, encode fine-grained semantics, and can handle intra-category object variations~\cite{manuelli2019kpam,adriaens2025spill}. Previously, obtaining keypoints required training task-specific detectors, limiting their scalability~\cite{adriaens2025spill,manuelli2019kpam,vecerik2021s3k}. Advances in vision foundation models~\cite{simeoni2025dinov3,heinrich2025radiov2.5,rombach2022stable-diffusion} now enable extracting keypoints using a single annotation and without task-specific training~\cite{amir2021ViTFeatures,tang2023dift}. Several works have leveraged this for keypoint-conditioned imitation learning~\cite{di2024KAT,wang2025skil,Levy2025P3PO,haldar2025pointpolicy}, demonstrating improved performance over raw RGB representations. However, many open questions remain, including how to best extract these keypoints, how to encode them for imitation learning and in which settings this approach outperforms other representations.

%what
This paper consolidates and builds on the methods proposed in previous works on keypoint-conditioned imitation learning (KIL)~\cite{di2024KAT,wang2025skil,Levy2025P3PO,haldar2025pointpolicy} and answers key questions about this approach. To do so, we evaluate the generalization capabilities, describe failure modes, investigate several design choices and expand the scope of KIL to tasks with multiple instances of the same object category.
Our pipeline for keypoint imitation learning consists of three parts. We start by manually specifying $3$-$6$ keypoints per object on a reference image. During training or inference, we use visual foundation models~\cite{bommasani2021foundation-models} to extract the corresponding keypoints across variations in poses, objects, and scenes. These keypoints are then lifted to 3D and used as input for a diffusion policy~\cite{chi2025diffusion_policy}. To make informed design choices, we compare different extraction methods (\textit{image matching}, \textit{instance matching} and \textit{tracking}), foundation models and model architectures.



\begin{figure}[t] 
    \centering 
    \includegraphics[width=0.99\linewidth]{imgs/first_page_success.png} \caption{Successful real-world rollouts of our keypoint imitation learning policies. The top-right label indicates the keypoint extraction method used: IM=image matching, IN=instance matching, T=tracking. (cf. Section~\ref{sec:method})} 
    \label{fig:first_page} 
    \vspace{-0.0cm}% <--- this magic squeezes the space between the caption and the rest of the text
    \vspace{-\baselineskip}
\end{figure}

%results
We perform over $2000$ real-world rollouts\footnote{\url{kil-manipulation.github.io}}, spanning two robot platforms and five different tasks, and evaluating both in-distribution performance, generalization to unseen objects, and scene variations. KIL achieves a $75\%$ success rate, compared to $47\%$ for an RGB diffusion %policy 
baseline, a gap that widens dramatically under scene variations (KIL $70\%$ vs.\ RGB $10\%$). KIL performs on par with S$^2$-diffusion~\cite{yang2025S2} ($73\%$), an object-centric baseline that uses estimated depth maps and segmentation masks. Different keypoint extraction methods and foundation models perform similarly, with performance varying from $73$\% to $85$\%. We also extend KIL to multi-instance tasks and illustrate how current foundation models struggle with large variations in object orientations.

Our contributions are summarized as follows:
\begin{itemize}

    \item A performant pipeline for KIL, based on an investigation of several design choices, including keypoint extraction methods, foundation models and policy architectures.
    \item A thorough assessment of its generalization capabilities and limitations, using over $2000$ real-world evaluations on two robot platforms.
    \item The first demonstration of KIL applied to tasks with multiple instances of the same object category, extending its scope.

\end{itemize}

\begin{figure*}[h]
    \centering
    \includegraphics[width=0.95\linewidth]{imgs/method.png}
    \caption{Overview of our three-stage keypoint imitation learning pipeline. 1) Keypoint references are manually annotated on a reference image for each object category. 2) Keypoints are extracted from each observation with visual foundation models and lifted to $3$D. 3) Keypoints are encoded by a transformer and used as input for a diffusion action head.}
    \label{fig:method-overview}
    \vspace{-0.0cm}% <--- this magic squeezes the space between the caption and the rest of the text
    \vspace{-\baselineskip}
\end{figure*}

\section{Related Work}
\label{sec:related-work}

\textbf{Keypoints for Robotic Manipulation} offer a compact, fine-grained and object-centric representation that has been used extensively~\cite{manuelli2019kpam,adriaens2025spill,vecerik2021s3k,chen2021unsupervised}.
Early works typically paired category-specific keypoint detectors with scripted policies~\cite{adriaens2025spill,manuelli2019kpam}. Others have focused on reducing labelling efforts~\cite{vecerik2021s3k,chen2021unsupervised}. More recently, visual foundation models~\cite{bommasani2021foundation-models} have enable 1-shot keypoint matching~\cite{amir2021ViTFeatures,tang2023dift,simeoni2025dinov3}, which can be used to build data-efficient robot manipulation systems~\cite{huang2025rekep,di2024KAT,wang2025skil}. %In this work
In this work, we focus on using such keypoint extraction pipelines for imitation learning.

% \begin{itemize}
%     \item kPAM - annotate in 3d
%     \item cloth keypoints - synthetic data for keypoints
%     \item S3K\cite{vecerik2021s3k} - clever consistency to reduce labeling effort
%     \item KETO - keypoints for tool manipulation with self-supervised discover/learning
%     \item unsupervised keypoint discovery - https://www.boyuan.space/unsup-3d-keypoints/
% other paper from MEL - TUSK.. not enough space but should add for dissetation.
% \end{itemize}

% Non-robot papers on keypoints for object representation ( outside of human poses..)
% \begin{itemize}
%     \item 
% \end{itemize}



\textbf{Object-Centric Imitation Learning} aims to improve the efficiency and generalization of sensorimotor imitation learning, which has progressed rapidly~\cite{barreiros2025LBM-1,chi2025diffusion_policy}, by using alternative representations to RGB images, often using visual foundation models~\cite{bommasani2021foundation-models} to avoid task-specific labeling. 

S$^2$-diffusion~\cite{yang2025S2} combines depth maps and segmentation images to provide a geometric and object-centric representation. GROOT~\cite{zhu2023learning} and DP3~\cite{ze20243d} construct pointcloud representations while VIOLA~\cite{zhu2023viola} derives object-centric proposals from pretrained models.
Rana et al.~\cite{rana2024learning} use oriented affordance frames and 6D object poses to boost generalization.
In this work, we use both standard diffusion policy~\cite{chi2025diffusion_policy} with RGB images and S$^2$-diffusion~\cite{yang2025S2} as baselines for KIL.
% \begin{itemize}
%     \item S2 diffusion
%     \item learning from 10 demos
%     \item DP3 -> point clouds
%     \item aira glasses -> object-centric, downsampled pointclouds
%     \item ... 
% \end{itemize}



\textbf{Visual Foundation Models for Keypoint Imitation Learning} have shown to improve in-distribution performance and generalization to unseen objects and scenes~\cite{Levy2025P3PO,di2024KAT,wang2025skil}. For keypoint specification, methods range from minimal human annotation on a single frame~\cite{haldar2025pointpolicy,Levy2025P3PO} to automatic distillation of task-driven points~\cite{vecerik2024robotap,zhang2025atk,wang2025skil,di2024KAT}. To extract keypoints from observations, these works either extract points from each frame via semantic descriptor matching~\cite{wang2025skil,di2024KAT} using various foundation models~\cite{tang2023dift,simeoni2025dinov3,heinrich2025radiov2.5}, or combine matching with explicit point-tracking~\cite{vecerik2024robotap,haldar2025pointpolicy,Levy2025P3PO,zhang2025atk}. 
%Some have also explored how keypoints, being embodiment-agnostic, can be used to learn directly from human observations~\cite{haldar2025pointpolicy,wang2025skil}. 
Finally, the control policies vary widely, encompassing visual servoing primitives~\cite{vecerik2024robotap}, supervised imitation networks~\cite{zhang2025atk,wang2025skil,Levy2025P3PO,haldar2025pointpolicy} and training-free, LLM-based in-context learning, which outperforms learned policies in ultra-low data regimes (e.g., $10$ demos) but loses this advantage as the number of demonstrations increases~\cite{di2024KAT}. 

In this paper, we consolidate these approaches and conduct an in-depth study of keypoint imitation learning, comparing a.o. keypoint extraction strategies, foundation models, and policy designs.

% \begin{itemize}
%     \item KAT - use keypoints as obs repr and action repr, LLM as planner with kp baseline.
%     \item SKIL - use keypoints as obs repr, train diffusion policy. use matching within episode. random references?
%     \item P3PO - use keypoints as obs repr, with BAKU. Use tracking within episode. - manual references?
%     \item Point Policy - use keypoints for cross-embodiment demonstration transfer
%     \item paper Mel - tracking with TAP for few-shot learning / replay.
%     \item ATK - keypoint selection for imitation learning https://www.arxiv.org/pdf/2506.13867v1
% \end{itemize}




\section{Method}
\label{sec:method}

Our method for keypoint imitation learning consists of three stages, illustrated in Figure~\ref{fig:method-overview}. First, a human operator specifies the semantic keypoints of interest by annotating a reference image for each object category (Section~\ref{sec:method-representation-specification}). Second, these keypoints are extracted from every observation, for both the demonstrations and during inference (Section~\ref{sec:method-representation-extraction}). Finally, the resulting sequence of $3$D keypoints is encoded by a transformer and used to condition a diffusion policy for action prediction (Section~\ref{sec:method-policy-learning}).

%  not sure if i want to add mathematical formalism for the process? only if it is more concise / clear to explain, because having math in a paper is not a goal by itself..


% $$P^i_{t} = f(I^i_t | I^i_{0..t-1}, \{(I,P)_{0...J}\})$$
% $$a_t = \pi(P_t||s_t) $$



\subsection{Specifying the Keypoint Representation}
\label{sec:method-representation-specification}
As in~\cite{Levy2025P3PO,haldar2025pointpolicy}, we manually annotate three to six $2$D keypoints on a reference image selected from a single demonstration, for each object category the robot interacts with. We follow three heuristics when selecting keypoints: minimize occlusions and spread keypoints across the object; place a keypoint near any task-critical contact point (e.g.,\ the heel of a shoe); and maintain a small distance from object borders. % to limit trailing artifacts in depth maps or inaccurate matches on background (image patches).

%\cite{di2024KAT,wang2025skil,zhang2025atk} have proposed methods to automate this one-time step of keypoint specification, yet we opted for manual specification due to the limited effort ($1-2$ minutes per task) involved. We believe there is value in finding those keypoints that optimize the matching accuracy and task performance, but consider this out of scope for this work.
\subsection{Representation Extraction}
\label{sec:method-representation-extraction}
Given the specification of the semantic keypoints, in the form of reference images with reference keypoints, we need to extract those keypoints across all initial states and for each timestep, both for training and inference.  We do this by first estimating the $2$D positions of the semantic keypoints and then lifting them to $3$D using the intrinsics and extrinsics of the calibrated RGB-D camera, as in~\cite{wang2025skil,Levy2025P3PO,di2024KAT}. % To lift the 2D keypoints to 3D, we use the depth map and intrinsics matrix to lift the pixels to 3D points and then express them in the frame of the robot, also used to express the actions and proprioceptive state, using the extrinsics matrix.

There are different ways to extract the $2$D positions of the semantic keypoints in each frame. In this work, we consider the following three fundamental approaches:

\textbf{Image Matching (IM)}
%This is the most straightforward approach, and 
Given the reference images and keypoints, we use a visual foundation model to compute a descriptor for each keypoint in the reference image. For each observation frame, we then compute these descriptors for every pixel in the image and use the cosine similarity with the reference descriptor of each keypoint to find the best match, as in~\cite{wang2025skil,di2024KAT}.


\textbf{Instance Matching (IN)}
This approach extends image matching by first running an open-vocabulary instance detector~\cite{carion2025sam3} (given a prompt for each category) to obtain $N$ instance masks per category, then restricting the descriptor-based cosine-similarity search to each object mask. 

This is a natural extension of the \textit{image matching} approach to handle multiple instances of the same category and potentially improve the matching accuracy. On the other hand, if the instance detector suffers from false positives or false negatives, this will impact the extracted keypoints significantly. An additional downside is the increased computational cost, as two foundation models must run for each observation. %\thomas{has been used for initial configuration in SKIL, and rekep?}


\textbf{Tracking (T)}
In the third and final approach, we use a dedicated tracking model to extract the keypoints, as in~\cite{Levy2025P3PO,haldar2025pointpolicy,vecerik2024robotap}.
To initialize the tracker, we use the \textit{instance matching} method to extract keypoints from the initial observation of a demonstration or rollout. Using separate models for tracking and initial matching can boost performance, especially with (short) object occlusions. 
%\thomas{should this be removed and only handled in the results section?}
%The main advantage of tracking over per-frame matching is that it can be more robust to partial occlusions, more precise due to temporal consistency, and allow for a computationally expensive initial frame matching approach while still being fast at inference time for all subsequent frames. It also achieves a cleaner separation of concerns between localization and tracking, which allows to use different datasets and training methods. However, the approach comes with several limitations. First, it adds engineering complexity, as the initial keypoints still need to be extracted via instance matching and the policy is no longer purely reactive. Second, all keypoints must be visible in the first frame for tracking to be initialized. Third, while the tracker can in principle handle brief occlusions, longer occlusions can lead to drift from which it cannot recover~\cite{wang2025skil}. \thomas{why can it not recover? could they not be smart about which frames to use in the rolling window?}


These three different approaches for extracting keypoints each have their (dis)advantages. We compare them extensively in our experiments and refer to Section~\ref{sec:experiments-extraction-comparison} for the results. In Section~\ref{sec:experiments-model-comparison}, we also compare different foundation models for each extraction method.


\subsection{Policy Learning}
\label{sec:method-policy-learning}

Given the success of diffusion policies for imitation learning~\cite{chi2025diffusion_policy}, we use a diffusion policy as action head to predict robot actions, combined with a transformer encoder for the keypoints, as in~\cite{wang2025skil}.

% \begin{itemize}
%     \item pre-norm
%     \item proprio as keypoints instead of concat?
%     \item handling multiple instances - padding with positional embeddings or identifying semantics?
%     \item history of obs from 4 to 2, to limit inference time..
%     \item 4 layers instead of 1.
%     \item ...
% \end{itemize}



The encoder closely resembles a standard transformer encoder~\cite{vaswani2017attention}, but we use pre-layer normalization, GeLU activations and no dropout.
%\thomas{refer to code appendix for details?} 
The keypoints and the proprioceptive state are each processed as separate tokens. To tokenize them, we project them to the transformer dimension with a learned layer and add sinusoidal positional embeddings to encode their semantics. The output tokens are mean-pooled into a single vector, following~\cite{wang2025skil}. We use a history of observations as input for the policy, each processed independently in this way; the resulting observation vectors are concatenated and used to condition the diffusion action head, which predicts the next chunk of actions.


The actions are encoded as $10$-dimensional vectors, representing the absolute EEF position, absolute rotation, and relative gripper state. The rotation is encoded as a six-dimensional vector representing the first two columns of the rotation matrix to make the actions continuous~\cite{chi2025diffusion_policy}.
The proprioceptive state is encoded as a $7$-dimensional vector representing the absolute position, absolute orientation using Euler angles, and the relative gripper state.
All inputs and actions are normalized to (-1,1) using min-max scaling. Unlike the original diffusion policy, we do not use Exponential Moving Averaging (EMA) during training.
Finally, we augment our data by adding zero-mean Gaussian noise to the proprioceptive state and keypoints and applying a random translation and rotation around the $z$-axis to the inputs and actions, as in~\cite{chisari2024flow-matching-pc-IL}. \appendixnote{We report the hyperparameters for the augmentations in Appendix~\ref{sec:appendix-implementation-details-and-hyperparameters} and have evaluated the impact of the augmentations on KIL performance in Appendix~\ref{sec:appendix-additional-experiments-augmentations}.}

\section{Experiments}
\label{sec:experiments}

With our experiments, we aim to answer the following questions about KIL.

\noindent \textbf{RQ-1} Does KIL boost generalization to unseen scenes and objects, compared to alternative methods? (Section~\ref{sec:experiments-generalization})

\noindent \textbf{RQ-2} How do the three keypoint extraction methods described in Section~\ref{sec:method-representation-extraction} perform? (Section~\ref{sec:experiments-extraction-comparison})

\noindent \textbf{RQ-3} Which visual foundation models are best suited for extracting the keypoint representations? (Section~\ref{sec:experiments-model-comparison})

\noindent \textbf{RQ-4} How should we encode the keypoints in the policy architecture? (Section~\ref{sec:experiments-encoder-comparison})

\noindent \textbf{RQ-5} Can we extend the use of keypoint imitation to tasks with multiple instances of the same object category? (Section~\ref{sec:experiments-multi-objects})

%Before detailing the results of these experiments, we first describe the hardware setup (Section~\ref{sec:experiments-hardware}), tasks (Section~\ref{sec:experiments-tasks}), baselines (Section~\ref{sec:experiments-baselines}) and experimental design (Section~\ref{sec:experiments-setup}). 

Next to answering these questions, we illustrate how current foundation models limit the performance of KIL under large object orientation variations in Section~\ref{sec:experiments-orientation}.



\subsection{Hardware Setup}
\label{sec:experiments-hardware}
Both hardware setups used in this work are shown in Figure~\ref{fig:robot_setup}.
The first setup uses a UR5e Robot with a Schunk EGK-$40$ gripper and a Zed2i RGB-D camera. We use a custom tactile fingertip~\cite{liu2025magtouch}, to limit the grasping force of the gripper.
The second setup consists of a Ufactory lite6 with an integrated gripper and an Orbbec Femto Bolt RGB-D camera.
%We use $2$ different sets of $3D$ printed grippers, adapted to grasp objects with different sizes, due to limitations in the gripper stroke width. 
In both setups, the scene cameras are placed as close as possible while having a view of the whole workspace. The cameras are opposing the robot at an inclination of approximately $45$ degrees, to limit occlusions by the gripper and self-occlusion of each object in the scene. A Meta Quest $3$ headset running Quest2ROS~\cite{welle2024quest2ros} is used for demonstration collection. %Both robots are also equipped with a wrist camera, but this camera is not used by any of the policies trained in this paper.

Policy and teleoperation actions are generated at $10$Hz and then interpolated to $50$Hz before sending them to the robot, to ensure smooth tracking. %For both setups, we use a foam pad to provide some passive compliance to the system, which proved important to somewhat avoid excessive forces being applied by the robot on the workspace/objects during policy rollouts. Finally, each robot has a workstation with an NVIDIA RTX$4090$ GPU that is used for training and inference.

\subsection{Tasks and Data Collection}
\label{sec:experiments-tasks}
To test the capabilities and limitations of KIL, we designed five tasks that require the robot to manipulate different object categories. The choice for these objects is largely motivated by prior work~\cite{Levy2025P3PO,wang2025skil,di2024KAT,huang2025rekep}. 
We collect $50$ demonstrations per task, taking about $1$ hour for each task.  This is slightly more than previous works~\cite{Levy2025P3PO,di2024KAT,wang2025skil}, but provides a practical trade-off between performance and data collection effort. We vary the position of the objects significantly during data collection, but, similar to~\cite{wang2025skil,di2024KAT}, we limit the variation in orientations to \(\pm 30^\circ\) around a canonical orientation (cf. Section~\ref{sec:experiments-orientation}). We collect demonstrations using two different objects, but without scene variations, as in~\cite{Levy2025P3PO,wang2025skil,di2024KAT}.
%Each 

The tasks and success criteria are described below. These criteria are combined with safety measures and a time limit (approx. twice the time it takes to teleoperate the task, excluding inference time) to determine if a rollout is successful.

\textbf{Place Shoe}: The goal is to pick a randomly placed shoe by the heel and put it on a tray that is placed at a fixed location. The task is successful if this happens without moving the tray significantly and with the shoe oriented within $90$ degrees of the y-axis of the tray.

\textbf{Place Mug} This task is similar to the previous one, but now requires picking a randomly placed, upright mug by the rim and placing it on a tray with a fixed position. 

\textbf{Pour Bottle}: The robot needs to pick a plastic 
bottle, move it over a cup, and tilt it as if to pour the contents into the cup.  The positions of the bottle and cup are randomized across the workspace. Since both objects are rotationally symmetric around the z-axis, the orientation does not matter for this task. The task is successful if the bottle is tilted beyond horizontal and is held directly above the cup. If the bottle is dropped, the task is terminated.

\textbf{Pick Pen}: The goal is to pick up a randomly placed pen and drop it in a cardboard box placed at a fixed location. The task is successful if the pen falls inside or remains on top of the box, without dropping the pen or moving the box significantly.

\textbf{Place Mouse}:  The goal is to grasp a randomly positioned mouse from the sides and place it on top of a mouse pad placed at a fixed location. The task is successful if the mouse is placed on the computer pad, facing the scene camera with the scroll wheel, without moving the pad significantly or dropping the mouse.

An example of each task and its keypoint representation can be seen in Figure~\ref{fig:first_page}, while the objects used for data collection and evaluation are shown in Figure~\ref{fig:robot_setup}. \appendixnote{The keypoint annotations (cf. Section~\ref{sec:method-representation-specification}) for each task are shown in Appendix~\ref{sec:appendix-keypoint-references}, while an overlay of the initial frames of the demonstrations for each task are shown in Appendix~\ref{sec:appendix-demo-configurations}.}

\begin{figure}[t] 
    \centering 
    \includegraphics[width=0.99\linewidth]{imgs/robot_setup.png} \caption{The $2$ robot setups and object sets used for each task (training objects grouped left of the evaluation objects for each task). } 
    \label{fig:robot_setup} 
    \vspace{-0.0cm}% <--- this magic squeezes the space between the caption and the rest of the text
    \vspace{-\baselineskip}
\end{figure}


\subsection{Baselines}
\label{sec:experiments-baselines}

We compare KIL against two baselines. The first is a standard RGB-based diffusion policy~\cite{chi2025diffusion_policy}. This uses the same diffusion action head as the keypoint policies, but now conditioned on the output of an image encoder that processes the RGB observation, resized to $320\times240$. We make one change to the encoder: we do not replace the global average pooling with a spatial softmax.
The second baseline is S$^2$-diffusion~\cite{yang2025S2}, which uses foundation models to generate a relative, normalized depth image and segmentation mask of the RGB image, using the policy architecture from~\cite{chi2025diffusion_policy}. %\thomas{discuss difference with keypoints.}

\begin{table*}[h!] % Use table* to span full width in 2-column papers
    \centering
    \renewcommand{\arraystretch}{1.2}
    \small % Reduced font size to fit 21 columns
    \setlength{\tabcolsep}{6pt} % Very tight padding to fit 21 columns
    \caption{Comparing keypoint extraction methods for KIL against baselines. IM=\textit{IMage matching}, IN=\textit{INstance matching},T=\textit{Tracking}. Each method is evaluated on $3$ sets of initial configurations, measuring in-distribution performance (IN), generalization to unseen objects (OBJ) and scene variations (S).}
    
    \begin{tabularx}{\textwidth}{l *{18}{Y} c c} 
        \toprule
        \multirow{2}{*}{\textbf{Method}} & 
        \multicolumn{3}{c}{\textbf{place shoe}} & 
        \multicolumn{3}{c}{\textbf{place mug}} & 
        \multicolumn{3}{c}{\textbf{pour cup}} & 
        \multicolumn{3}{c}{\textbf{pick pen}} & 
        \multicolumn{3}{c}{\textbf{place mouse}} & 
        \multicolumn{3}{c}{\textbf{ Task Avg. (\%)}} &
        \multirow{2}{*}{\textbf{Avg. (\%)}} &
        \multirow{2}{*}{\textbf{CI (\%)}} \\
        
        \cmidrule(lr){2-4} \cmidrule(lr){5-7} \cmidrule(lr){8-10} \cmidrule(lr){11-13} \cmidrule(lr){14-16} \cmidrule(lr){17-19}
        
        & \textbf{in} & \textbf{obj} & \textbf{s} & \textbf{in} & \textbf{obj} & \textbf{s} & \textbf{in} & \textbf{obj} & \textbf{s} & \textbf{in} & \textbf{obj} & \textbf{s} & \textbf{in} & \textbf{obj} & \textbf{s} & \textbf{in} & \textbf{obj} & \textbf{s} & & \\
        \midrule
        RGB~\cite{chi2025diffusion_policy}           & 8  & 6 & 3 & 10 & 7 & 0  & 4 & 0 & 1 & 7 & 7 & 0 & 8  & 8 & 1 & 74\% & 56\% & 10\% & 47\% & {\scriptsize [38, 55]} \\
        S$^2$~\cite{yang2025S2}            & 9  & 9 & 5 & 8  & 9 & 10 & 6 & 3 & 5 & 7 & 7 & 5 & 10 & 9 & 8 & 80\% & 74\% & 66\% & 73\% & {\scriptsize [66, 80]} \\
        KIL (IM)            & 10 & 6 & 9 & 10 & 8 & 8  & 7 & 4 & 5 & 7 & 8 & 4 & 10 & 8 & 9 & 88\% & 68\% & 70\% & 75\% & {\scriptsize [68, 82]} \\
        KIL (IN) & 10 & 6 & 5 & 10 & 7 & 10 & 1 & 4 & 2 & 8 & 9 & 6 & 9  & 9 & 9 & 76\% & 70\% & 64\% & 70\% & {\scriptsize [62, 77]} \\
        KIL (T) & 9  & 6 & 6 & 9  & 9 & 9  & 8 & 8 & 8 & 6 & 8 & 7 & 6  & 6 & 7 & 76\% & 74\% & 74\% & 75\% & {\scriptsize [67, 81]} \\
        \bottomrule
    \end{tabularx}
    \label{table:single-object-main}
    \vspace{-0.0cm}% <--- this magic squeezes the space between the caption and the rest of the text
    \vspace{-\baselineskip}
\end{table*}

\subsection{Experimental Setup}
\label{sec:experiments-setup}
\textbf{Implementation Details}
For \textit{image matching (IM)}, we use RADIOv2.5-B~\cite{heinrich2025radiov2.5} as feature extractor. For \textit{instance matching (IN)}, we use SAM3~\cite{carion2025sam3} to predict instance masks and DIFT~\cite{tang2023dift} with Stable Diffusion 2.1~\cite{rombach2022stable-diffusion} as feature extractor.
%which outperformed RADIOv2.5-B in this setting (cf.\ Section~\ref{sec:experiments-model-comparison}), although being significantly slower
We use an ensemble size of $2$ for DIFT to balance inference cost and performance. For \textit{tracking (T)}, we use CoTracker3~\cite{karaev2025cotracker3} with the online modification from~\cite{Levy2025P3PO}; the initial frame is matched using SAM3 and DIFT with an ensemble size of $8$ to reduce variance due to the random noise initialization of DIFT. All images are resized to $640\times480$. We use bilinear interpolation to upscale the output of all feature extractors to the original size, as in~\cite{wang2025skil,huang2025rekep}.

Our keypoint transformer encoder consists of $4$ layers, using a token size of $128$ and an MLP size of $512$. The number of learnable parameters is approximately $1M$. As action head, we use the Unet from~\cite{chi2025diffusion_policy} and keep most hyperparameters at their default values. The observation history is set to $2$.
%We believe covering a larger period of time in the observation history might be helpful with occlusions, but it also increases the computational cost of the policy.
For all policies, we predict $16$ actions and execute the next $8$ actions before predicting the next action chunk.

For keypoint-based policies, we halve the number of channels in the Unet, since we did not observe improvements by using a larger action head, in line with~\cite{lin2025imitation_scaling_laws}. We also use a larger batch size of  $128$, compared to $32$ for the image-based baselines. All policies use a learning rate of $1e^{-4}$ with a cosine scheduler, as in~\cite{chi2025diffusion_policy}.
We train image-based baselines for $100k$ steps, and KIL for $250k$ steps. %We found it beneficial to train the keypoint-based policies longer, especially when using augmentations. This is similar to how state-based policies are trained for more steps in~\cite{chi2025diffusion_policy}.
\appendixnote{Additional details on the implementation, hyperparameters and inference latencies are reported in Appendix~\ref{sec:appendix-implementation-details-and-hyperparameters}.}

\textbf{Evaluation protocol} 
We perform $30$ rollouts per task: $10$ with in-distribution initial configurations, $10$ with unseen objects (using $3-5$ objects per task, shown in Figure~\ref{fig:robot_setup}), and $10$ with scene variations, which include distractors and changing the background color of the surface. %Unlike~\cite{}, we do not test for out-of-distribution object poses because we believe it is feasible to cover the poses you are interested in during data collection.
We use the same initial scene configurations \appendixnote{(cf. Appendix~\ref{sec:appendix-eval-configurations})} for the different policies to ensure fairness by overlaying reference images on the current scene and manually positioning all objects accordingly. Additionally, we control lighting conditions and background elements to limit distribution shifts between evaluations.

for the different policies to ensure fairness by overlaying reference images on the current scene and manually positioning all objects accordingly. Additionally, we control lighting conditions and background elements to limit distribution shifts between evaluations.

\textbf{Statistical Analysis} 
Next to providing point estimates of the aggregated success rate or task completion, we provide 95\% confidence intervals (CI) to quantify uncertainty and assess statistical significance. For binary success results, we assume a binomial distribution, as in ~\cite{barreiros2025LBM-1} 
and report Clopper-Pearson CIs for the average success rate. For task completion scores, we
rely on the central limit theorem ($N=60$) to assume a normal distribution of the average completion rate, as in~\cite{barreiros2025LBM-1}. We report the standard Wald CI $\hat{\mu} \pm 1.96\,\hat{\sigma}/\sqrt{N}$. \appendixnote{Finally, for tables that only include results aggregated over multiple tasks or evaluation scenarios, we report the full results in Appendix~\ref{sec:appendix-detailed-results}.}


%we follow the methodology of~\cite{barreiros2025LBM-1} to assess the statistical significance of the results. For each experiment, we use pairwise testing for all methods under comparison, using Bonferroni correction, and report which results differ significantly using CLD: Methods which do not share the same letter are significantly different at the $5\%$ level. For binary policy success rates, we assume a binomial distribution and use Fisher's exact test to determine if the difference in average success rate is significant. For the task completion metrics, we use a Welch t-test on the average task completion rates, as in~\cite{barreiros2025LBM-1}. The use of the Fisher test for binary success rates differs from \cite{barreiros2025LBM-1}, where a sequential test is used with varying sample sizes. We did not use varying sample sizes and hence opted for the exact Fisher test.
%\thomas{we should maybe complement p-values with CI. If that is the case, not sure the tests contribute much actually, we could just use a non-overlapping CI to determine significance?}

\subsection{Comparing KIL Against Baselines}
\label{sec:experiments-generalization}
In this first experiment, we compare the performance of KIL against the two baselines, standard RGB-based diffusion and S$^2$ (\textbf{RQ-1}), across all five tasks described in Section~\ref{sec:experiments-tasks}.
The results are presented in Table~\ref{table:single-object-main}. KIL performs significantly better than the RGB-based diffusion policy overall, with the RGB policy achieving only $47\%$ success compared to $75\%$ for KIL with \textit{image matching}, the best-performing extraction method (cf.\ Section~\ref{sec:experiments-extraction-comparison}). While the RGB policy matches KIL on in-distribution evaluations and generalizes surprisingly well to unseen objects with varying appearance and geometry, it achieves only $10\%$ success on scene variations. S$^2$, by contrast, performs on par with KIL even under scene variations, which noticeably alter the depth maps relative to the demonstrations. Compared to prior work~\cite{di2024KAT,Levy2025P3PO,wang2025skil}, KIL achieves similar task performance, while our RGB policies perform slightly better, which we attribute to the larger number of training demonstrations.
%\thomas{relate to KAT/p3po/skil?  kp performace is similar to skil and KAT. RGB is 40 for KAT(40 demos), but a bit smaller for SKIL., but we use more demos and also have better implementation?}

Imprecise grasping was the most prominent failure mode across all policies: the robot nearly misses the object multiple times until a timeout occurs, or obtains a poor grasp that causes the rollout to fail later on, e.g.\ by dropping the object in a way that recovery is not possible. For KIL specifically, grasp failures occurred even when the extracted keypoints appeared accurate, suggesting insufficient policy generalization. Even in successful rollouts, KIL occasionally exhibited reduced grasping precision, which we attribute to variance and inaccuracies in the extracted keypoints.
%The RGB baseline is particularly susceptible to scene distractors: when distractors are present or the background colour changes, the policy frequently attempts to grasp distractor objects rather than the target, hence the very low succes rate for scene generalization.

%\thomas{something about getting stuck? was this only for me or also for Marco?}

%Finally, some policies occasionally command large, discontinuous end-effector motions that exceed the safety limits of the robot controller and do not recover from this, causing termination of the rollout.

\subsection{Comparing Keypoint Extraction Methods}
\label{sec:experiments-extraction-comparison}
We now compare the three methods for extracting keypoint representations discussed in Section~\ref{sec:method-representation-extraction}: \textit{image matching (IM)}, \textit{instance matching (IN)} and \textit{tracking (T)} (\textbf{RQ-2}).
In Table~\ref{table:single-object-main}, we find that all three methods perform similarly, achieving success rates of $75\%$, $70\%$ and $75\%$, averaged over the five tasks and different evaluation settings.

Empirically, \textit{matching}-based approaches for keypoint extraction suffer more from occlusions than the \textit{tracking} approach. At the same time, \textit{tracking} is more prone to drifting over time if objects are occluded for too long, or when the extracted keypoints of the first frame are inaccurate. 
%Matching-based approaches, being stateless, do not suffer these issues although they have more spurious matches. 
We illustrate both keypoint failure modes in Figure~\ref{fig:failures}. 

Since we found no noticeable differences in performance between the extraction methods, we have used \textit{image matching}, which is simpler and computationally less expensive, for all experiments in this paper, unless stated otherwise.



\subsection{Comparing Feature Models for Keypoint Matching}
\label{sec:experiments-model-comparison}

To find the most suitable keypoint feature model for KIL (\textbf{RQ-3}), we compare three foundation models across different keypoint extraction methods on two tasks: \textit{Place Shoe} and \textit{Place Mug}. The three models are RADIOv2.5-B~\cite{heinrich2025radiov2.5}, DIFT~\cite{tang2023dift} and DINOv3-B~\cite{simeoni2025dinov3}.

%results
Table~\ref{tab:model-comparison} shows that RADIOv2.5-B ($85\%$) performs better than DINOv3-B ($73\%$) and DIFT ($62\%$) for \textit{image matching}. We noticed that, although DIFT produces very accurate matches without occlusions, it is more prone to having incorrect matches, especially when the actual object is partially occluded and with scene variations. Using either \textit{instance matching} or \textit{tracking} reduces the impact of this issue, as can be seen in Table~\ref{tab:model-comparison}.
Finally, for \textit{instance matching}, and \textit{tracking}, we find no significant differences between RADIOv2.5-B and the DIFT models with ensemble sizes of $2$ and $8$. 

Based on these results, we have used RADIOv2.5-B for \textit{image matching}, DIFT$_2$ for \textit{instance matching} and DIFT$_8$ for \textit{tracking} in this paper, unless stated differently.
%As already discussed in Section~\ref{sec:experiments-single-object-main}, there are  no significant differences between the different keypoint extraction methods. 


%\cite{wang2025skil} also compared different foundation models for imitation learning with image matching. They also found radio models to outperform Dino, but in their case DIFT slightly outperformed radio. We compared their implementation of DIFT against ours, but could not find anything that explains this difference. We hypothesize that their scenes contain less visual noise compared to ours (which has a very cluttered background, cf Figure~\ref{TODO} \thomas{refer to a figure in which my setup background is visible}), limiting the potential of spurious matches.

%\thomas{KAT uses another method. tried it but 1) it is super expensive and 2) does not outperform modern ViTs. 3) does not work well for modern models ( larger patch size?)}

%relate to 2D metrics.
%only if we have space..


\begin{table}[h]
    \centering
    \renewcommand{\arraystretch}{1.2}
    % \small % Slightly reduces font size to help with fit
    \setlength{\tabcolsep}{5pt} % Reduces the padding between columns
    \caption{Comparison of keypoint feature models for the different keypoint extraction methods.}
    \begin{tabularx}{\linewidth}{l l Y Y c c}
        \toprule
        \makecell[l]{\textbf{Extraction} \\ \textbf{Method}} & \makecell[l]{\textbf{Feature} \\ \textbf{Model}} & \makecell[c]{\textbf{place} \\ \textbf{shoe} \\ \textbf{(/30)}} & \makecell[c]{\textbf{place} \\ \textbf{mug} \\ \textbf{(/30)}} & \makecell{\textbf{Avg.} \\ \textbf{(\%)}} & \makecell{\textbf{CI} \\ \textbf{(\%)}} \\
        \midrule
        \multirow{3}{*}{\textit{Image}} & \textbf{RADIOv2.5-B}  & 25 & 26 & 85\% & {\scriptsize [73, 93]} \\
        % & RGB & 17 & 17 & 57\% & \\
        & DIFT$_2$  & 13 & 24 & 62\%  & {\scriptsize [48, 74]} \\
        & DINOv3-B & 20 & 24 & 73\%  & {\scriptsize [60, 84]} \\
        \midrule
        \multirow{2}{*}{\textit{Instance}} & RADIOv2.5-B & 19 & 27 & 77\% & {\scriptsize [64, 87]} \\
        & \textbf{DIFT$_2$} & 21 & 27 & 80\% & {\scriptsize [68, 89]} \\
        \midrule
        \multirow{3}{*}{\textit{Tracking}} &RADIOv2.5-B & 20 & 24 & 73\% & {\scriptsize [60, 84]} \\
        & DIFT$_2$ & 23 & 26 & 82\% & {\scriptsize [70, 90]} \\
        & \textbf{DIFT$_8$} & 21 & 27 & 80\% & {\scriptsize [68, 89]} \\
        \bottomrule
    \end{tabularx}
\label{tab:model-comparison}
\end{table}


\subsection{Comparing Keypoint Encoders}
\label{sec:experiments-encoder-comparison}

In this section, we aim to find the most suitable way to encode the keypoints before passing them to the diffusion action head (\textbf{RQ-4}). We compare our transformer encoder, described in Section~\ref{sec:method-representation-extraction}, against two alternatives. The first forgoes any encoder and simply concatenates the keypoints to the proprioceptive state of the diffusion policy, as in~\cite{di2024KAT}. The second uses the same transformer encoder but augments each keypoint's $3$D position with the cosine similarity against all keypoint feature descriptors, as in~\cite{wang2025skil}.

%results
The results are shown in Table~\ref{tab:encoder-comparison}. We find that using no encoder performs on par with our transformer encoder, both achieving $85\%$ success rate. Appending cosine similarities to the keypoint $3$D positions, however, degrades performance dramatically to $7\%$. We hypothesize that, in combination with the geometric augmentations used in our pipeline, the network overfits on the similarity scores, which carry spurious information that is easier to memorize than the augmented $3$D positions. Removing augmentations substantially recovers performance for the similarity-based variant to $65\%$, but this is still below the performance we obtained without similarity scores. Beyond this negative interaction with augmentations, we also found the cosine similarity scores to be poorly calibrated across keypoints and scenes. Based on the results from this experiment, we have used the transformer encoder without similarities for all experiments in this paper, unless stated otherwise.
\begin{table}
    \centering
    \renewcommand{\arraystretch}{1.2}
    \setlength{\tabcolsep}{5pt}
    \caption{Encoder Comparison. Evaluated on the \textit{Place Shoe} and \textit{Place Mug} tasks.}
    \begin{tabularx}{0.6\linewidth}{l Y c}
        \toprule
        \textbf{Encoder} & \textbf{Avg.} & \textbf{CI (\%)} \\
        \midrule
        None              & 85\% & {\scriptsize [73, 93]} \\
        \textbf{ours}     & 85\% & {\scriptsize [73, 93]} \\
          \hspace{.3em} + sim & 7\% & {\scriptsize [2, 16]} \\
            \hspace{1.2em} - aug & 65\% & {\scriptsize [52, 77]} \\
        \bottomrule
    \end{tabularx}
    \label{tab:encoder-comparison}
    \vspace{-0.0cm}% <--- this magic squeezes the space between the caption and the rest of the text
    \vspace{-\baselineskip}
\end{table}

% \subsection{Comparing instance models}
% \thomasRemark{not sure if we will have time to do this!}

% Which instance model to use?

% trade-off between Performance, inference time..

% \begin{itemize}
%     \item SAM3
%     \item gDINO
%     \item OWL-ViT
% \end{itemize}


\subsection{Multi-Object Tasks}
\label{sec:experiments-multi-objects}
% \begin{table}[h] % Use table* to span full width in 2-column papers
%     \centering
%     \renewcommand{\arraystretch}{1.2}
%     \small % Reduced font size is highly recommended for 14 columns
%     \setlength{\tabcolsep}{5pt} % Very tight padding to fit 14 columns
%     \caption{multiple objects comparison}
    
%     % defined as: Method column (l) + 12 dynamic columns (Y) + Average column (C)
%     \begin{tabularx}{\textwidth}{l *{12}{Y} C} 
%         \toprule
%         \multirow{2}{*}{\textbf{Method}} & 
%         \multicolumn{3}{c}{\textbf{place 2 shoes}} & 
%         \multicolumn{3}{c}{\textbf{place 2 mugs}} & 
%         \multicolumn{3}{c}{\textbf{pick 2 pens}} & 
%         \multicolumn{3}{c}{\textbf{place 2 mice}} & 
%         \multirow{2}{*}{\textbf{Avg. ($CI_{95\%}$)}} \\
        
%         \cmidrule(lr){2-4} \cmidrule(lr){5-7} \cmidrule(lr){8-10} \cmidrule(lr){11-13}
        
%         & in & obj & scene & in & obj & scene & in & obj & scene & in & obj & scene & \\
%         \midrule
%         \# rollouts & /10 & /10 & /10 & /10 & /10 & /10 & /10 & /10 & /10 & /10 & /10 & /10 & /120 \\
%         RGB &  6.5& 4.5 &1.5  & 8 & 4 & 2 & 8.5 & 4.5 & 2.5 & 8 & 9 & 3 & 62\\
%         S2 &  &  &  &  &  &  &  &  &  &  &  &  & \\
%         KP + instance &6  & 1.5 & 5 & 6 & 2.5&  4&  6.5 & 3.5 & 6.5 & 8 & 5 & 5  & 59.5 \\
%         KP + tracking  &  &  &  &  &  &  &  &  &  &  &  &  & \\
%         \bottomrule
%     \end{tabularx}
% \end{table*}
A limitation of using \textit{image matching} with multiple objects of the same category is that the keypoints might not be matched on each object, as multiple pixels with the highest similarity scores to a reference keypoint could be located on the same object, thus resulting in incomplete keypoint representations.
In this section, we extend KIL to such tasks (\textbf{RQ-5}) by extending the tasks from Section~\ref{sec:experiments-tasks} to make the robot handle two objects, instead of one.
During data collection, we have the robot handle the object closest to the camera first, to limit occlusions.

Model architecture, parameters, and keypoint references remain the same, with three exceptions:
(i) To account for the increased length of the tasks, we increase the number of training steps to $150K$ for all image-based baselines and $350K$ for KIL. % to keep the number of epochs in the same order of magnitude as before. 
(ii) Some changes are made to the extraction methods: For \textit{image matching}, we add a second reference image and keypoints to the configuration, allowing to extract keypoints for up to two objects. For \textit{instance matching}, we extract the two instances with the highest confidence and match keypoints on both masks. Should there be only one instance, we pad the list of keypoints with zeros, so that the total number of input tokens for the transformer encoder is consistent.
(iii) The token pooling strategy in the keypoint encoder is changed from mean-pooling to the output of a newly-added class token, concatenated with a maxpool over all other tokens, as in PointBERT~\cite{yu2022point-bert}. We made this change because the initial performance for the two-object tasks using mean-pooling was very poor, even though the keypoints were reasonably accurate. For the tasks from Section~\ref{sec:experiments-tasks}, changing the pooling strategy did not result in significant differences. \appendixnote{Our experiment with different pooling strategies is reported in Appendix~\ref{sec:appendix-additional-experiments-token-pooling}.}


\begin{table}
    \centering
    \renewcommand{\arraystretch}{1.2}
    \small 
    \setlength{\tabcolsep}{4pt}
    \caption{Task completion for tasks in which multiple instances of the same object category are handled. }
    
    \begin{tabularx}{0.95 \linewidth}{l Y Y c c}
        \toprule
        \textbf{Method} & \textbf{place 2 shoes (/30)} & \textbf{pick 2 pens (/30)} & \textbf{Avg. (\%)} & \textbf{CI (\%)} \\
        \midrule
        RGB~\cite{chi2025diffusion_policy}     & 13   & 15.5 & 48\% & {\scriptsize [37, 58]} \\
        S$^2$~\cite{yang2025S2}      & 14.5 & 20.5 & 58\% & {\scriptsize [47, 69]} \\
        KIL (IM) & 21.5 & 16.5 & 63\% & {\scriptsize [54, 73]} \\
        KIL (IN) & 21.5 & 15   & 61\% & {\scriptsize [50, 72]} \\
        KIL (T)  & 15.5 & 15.5 & 52\% & {\scriptsize [40, 63]} \\
        \bottomrule
    \end{tabularx}
    \label{table:2-object-experiments}
    \vspace{-0.0cm}% <--- this magic squeezes the space between the caption and the rest of the text
    \vspace{-\baselineskip}
    % \vspace{-\baselineskip}
\end{table}

% results: KP > RGB, as for single object. KP similar to S2 as for 1 object. 

The results of our experiments are shown in Table~\ref{table:2-object-experiments}. Note that instead of the final success rate, we now report task completion (defined as $0.5$ per correctly handled object) to better reflect model capabilities on these longer-horizon tasks~\cite{barreiros2025LBM-1}.  We find that KIL, irrespective of the extraction method, outperforms RGB ($48\%$) and performs similarly to S$^2$ ($58\%$). These results are consistent with the findings in Section~\ref{sec:experiments-generalization}, although the difference between RGB and KIL is smaller.

% The performance of \textit{image matching (IM)}  is surprising, given the concern about incomplete keypoint representations. We noticed that this problem does indeed happen frequently, but observed that the keypoints are mostly matched on the most prominent object in each frame. By coincidence, this aligns with the data collection strategy and task specifics, and hence there are always sufficient keypoints on the objects the robot is handling. From this, we conclude that for our tasks, image-based matching can work, but in general, we expect it to break when the object being manipulated is not the most prominent one in the image.

% would love to illustrate this, but there is no space to do so.. 
The strong performance of \textit{image matching (IM)} is surprising given the risk of incomplete representations. In practice, keypoints tended to match onto the most prominent object, which coincidentally aligned with the data collection strategy used in our experiments. We expect \textit{image matching} to break when the manipulated object is not the most visually prominent.

Next to the failure modes already described in Section~\ref{sec:experiments-extraction-comparison}, we noticed two additional failure modes in the multiple object experiments.
Firstly, for \textit{instance matching (IN)}, irrelevant scene elements can be falsely detected as a second object, causing failures if this did not happen during training. This is illustrated in Figure~\ref{fig:failures}.  Tuning the confidence threshold to avoid these false positives is non-trivial since a partially occluded true positive may score lower than a distractor.
Secondly, with \textit{tracking}, the tracker can snap onto another object of the same category when they move in front of each other, causing the policy to not see the second object. Per-frame matching approaches recover from this immediately once the second object is visible again.


% discuss validate hypothesis for token pooling.
% also try to argue why it does not matter for the other experiments with 1 object (it would impact all keypoint models equally, and only by a small amount) 
% To conclude this section, we elaborate on the use of a different token aggregation strategy. As stated before, we found that using pointBERT token aggregation (CLS token + maxpool) worked better than using mean-pooling, which we adopted from~\cite{wang2025skil}.
% We hypothesize that as the number of (independent) keypoints (and thus tokens) increases, the mean-pooling reduces spatial accuracy by blurring the relevant values. In Table~\ref{table:token-pooling-comparison}, we compare mean pooling against pointBERT pooling. We also compare against widening the transformer encoder, thus providing more capacity to the tokens, or not using an encoder at all. We find that pointBERT  outperforms all other options for the two-object tasks.
% We expect that the performance for the single-object tasks would also improve slightly with pointBERT pooling. This is illustrated by the comparison on the \textit{pour} task in Table~\ref{table:token-pooling-comparison}. However, we expect it to slightly boost performance for all keypoint-based methods equally, hence not affecting the conclusions of previous experiments.

\subsection{Robustness Against Rotations}
\label{sec:experiments-orientation}

% Less obvious limitations are the limited ability of current matching models to handle orientation changes, and the complete lack of calibration of the matching models, which makes them snap keypoints on another part of the image, if the actual semantic correspodence is not visible. This inability of foundation models to handle orientations well made us limit the orientations of the objects in our tasks, similar to others~\cite{Levy2025P3PO,di2024KAT,wang2025skil}.

A surprising limitation of current foundation models used for keypoint features is their poor robustness to large object orientation changes, together with their lack of calibration. Because of this, we restricted object orientations for all tasks in this paper, consistent with prior work~\cite{Levy2025P3PO,di2024KAT,wang2025skil}. To better illustrate this limitation, we now create task variants in which object poses span the full \(\pm 180^\circ\), for both the \textit{Place Mouse} and \textit{Pick Pen} tasks. We also impose an orientation constraint on the task success criteria: an execution is counted as successful only if the gripper grasps the object with an orientation that is within \(\pm 30^\circ\) of the orientation used consistently during teleoperation. 
\begin{table}
    \centering
    \renewcommand{\arraystretch}{1.2}
    % \small % Slightly reduces font size to help with fit
    \setlength{\tabcolsep}{5pt} % Reduces the padding between columns
    \caption{Illustration of performance drop when increasing the range of initial orientations for task objects. Evaluated on the \textit{Place Mouse} and \textit{Pick Pen} tasks.}
    \begin{tabularx}{0.95 \linewidth}{l Y Y c Y}
        \toprule
        & \multicolumn{2}{c}{\(\boldsymbol{\pm 180^\circ}\)} & \(\boldsymbol{\pm 30^\circ}\) & \(\boldsymbol{\Delta}\) \textbf{(30$\to$180)} \\
        \cmidrule(lr){2-3} \cmidrule(lr){4-4}
        \textbf{Method} & \textbf{Avg. (NS)} & \textbf{Avg.} & \textbf{Avg.} & \\
        \midrule
        RGB      & 70\% & 47\% & 52\% & $-10\%$ \\
        S$^2$    & 80\% & 70\% & 77\% & $-9\%$ \\
        KIL (IM) & 45\% & 35\% & 77\% & $-55\%$ \\
        KIL (T)  & 48\% & 48\% & 67\% & $-28\%$ \\
        \bottomrule
    \end{tabularx}
\label{tab:more-rotations-experiment}
\vspace{-0.0cm}% <--- this magic squeezes the space between the caption and the rest of the text
    \vspace{-\baselineskip}
\end{table}


Table~\ref{tab:more-rotations-experiment} reports the aggregated performance across all initial configurations for both tasks (Avg). To better expose the degradation of KIL under larger orientation variations, we also report an aggregate that excludes scene variation evaluations (\textit{Avg. (NS)} - no scene), since RGB performs poorly in this setting, which can partially mask the gap. We also report the performance on the corresponding \(\pm 30^\circ\) tasks, in addition to the relative difference between both ( $\Delta$ (30$\to$180)).
The results show a clear trend: keypoint-based methods degrade substantially when object orientation spans the full \(\pm 180^\circ\) range and correct gripper approach orientation is required, with KIL (IM) dropping by $55\%$ and KIL (T) by $28\%$. In contrast, RGB and S$^2$ are affected considerably less, dropping by only $10\%$ and $9\%$ respectively.

\begin{figure*}
    \centering
    \includegraphics[width=0.99\linewidth]{imgs/failures_cut.png}
    \caption{Failure examples for the keypoint extraction methods. Left (\textit{image matching}): matching assigns keypoints to incorrect parts of the scene rather than the intended object regions; the reference keypoints are shown in the bottom-right. Center (\textit{tracking}): drift causes some keypoints to snap onto a different object. Right (\textit{instance matching}): a false positive instance detection introduces additional keypoints. }
    \label{fig:failures}
    \vspace{-0.0cm}% <--- this magic squeezes the space between the caption and the rest of the text
    \vspace{-\baselineskip}
\end{figure*}

\section{Discussion and Limitations}

Similar to~\cite{di2024KAT,wang2025skil,Levy2025P3PO}, we found KIL to outperform RGB-based policies. However, as discussed in Section~\ref{sec:experiments-generalization}, we found the difference to be most significant under scene variations. 
%Based on this observation
Therefore, an interesting direction for future work is to compare KIL against methods that augment the demonstrations with image generation models~\cite{yu2023ROSIE} and test if the difference between KIL and RGB-based policies persists.

Unlike \cite{wang2025skil,Levy2025P3PO}, we found our object-centric baseline, S$^2$-diffusion, to perform on par with KIL. This raises the question of why the semantic granularity and compactness of keypoints does not translate into superior performance over representations such as masks or depth maps. A possible explanation is that the extraction pipelines (i.e. the foundation models) are more performant for those other representations.

We observed KIL to be less precise and degrade faster under occlusions and variations in object orientation than RGB-based policies. These are limitations of the current foundation models used to extract keypoints, rather than the concept of using keypoints for imitation learning. Potential directions to overcome them include end-to-end finetuning of the extraction pipeline or combining keypoints with other representations~\cite{chen2025history-keypoint-IL,cao2025COIL}, to combine their strengths.

%\thomas{mention that depth and 3d lifting were not observed to cause failures? depth not perfect, but seems good enough? but are we sure of this? hard to know actually, but every time I looked at robot EEF in rerun, its relative position wrt the 3d keypoints seemed ok compared to the actual scene. in \cite{haldar2025pointpolicy} they found depth to not work at all, and resorted to multi-view triangulation. feel like the use of realse d435 is the culprit here, that one has shitty depth. (cf DP3, where they use the L series and explicitly state that D400 wont work)}

%A final interesting direction for future work is to extend the use of keypoints for human-robot embodiment mapping~\cite{haldar2025pointpolicy} to cross-robot generalization and sim-to-real transfer of keypoint-based policies.

% robot that can perform the task, should be able to point to a semantic location. question is if you need to therefore limit the observation to these points? what if you simply use this as addition training object somehow, but then still give it the full image? this is closely related to how VLMs are used as starting point for VLAs, and are even finetuned specifically for robotics initialy (Gemini-Robotics-ER). 

We also identify several limitations of our work.
We did not consider the use of wrist cameras, which help to reduce occlusion and can boost generalization for RGB-based policies~\cite{hsu2022wrist-cam-policy-benefits}. We also used depth sensors in this work. Although we found their precision to be sufficient, the use of depth sensors limits the scope of tasks to a.o. objects that are not reflective or transparent. % (no transparent liquids in the bottles).
% Interestingly, we obtained good performance with depth sensors for the tasks in this paper, and did not need to resort to multi-view triangulation as in~\cite{haldar2025pointpolicy}.
% Triangulation from multiple views~\cite{haldar2025pointpolicy} or the use of learned (non-metric) depth estimators~\cite{Levy2025P3PO,yang2025S2} provide alternative solutions. 
The performance of SAM3~\cite{carion2025sam3} for \textit{instance matching} depends a.o. on how discriminative the prompt is. For example, in the \textit{Pick Pen} task, we used "paper box" instead of "box" as prompt to improve performance. Such prompt engineering increases the time required for representation specification.
%Although we performed extensive experiments with different design choices, many hyperparameters remain underexplored and could be optimized for better performance.
Finally, our investigation of KIL's generalization capabilities is limited to $2$-shot object generalization and $1$-shot scene generalization. A more extensive evaluation of how generalization scales with the number of demonstrated variations, as in~\cite{lin2025imitation_scaling_laws}, would provide a more complete picture. 

%We believe simulation would help to make this feasible and ensure identical data distributions for training and evaluation. \thomas{can also have sim-to-real gap..}


% \begin{itemize}
%     \item what does it mean that S2=keypoints? is keypoint repr not better than simple masks + depth? or is it that extraction pipelines are more consistent/better for S2?
%     \item RGB < others mostly because of scene distractors, could try to solve this with data augmentations?
% \end{itemize}
% Limitations 
% \begin{itemize}    
%     \item Occlusions !!!
%     \item object categories that are not well supported by the foundation models.
%     \item limited precision, should aim to finetune foundation models? and need better baselines? (cf comparison).
%     \item rotations heavily affect keypoint matching.
%     \item using wrist cameras is non-trivial (depth might not be feasible since it does not work when the end effector is close to the table).
%     \item inference latency. sam3 and S2 slower due to segmentation pass.
%     \item light reflections for depth can be catastrophic, anydepth as possible solution.
%     \item large number of hyperparamters are not validated. but this is a problem for many robot learning methods. 
% \end{itemize}

% Future work
% \begin{itemize}
%     \item finding the best keypoint configurations automatically. 
%     \item other tracking methods?
%     \item finetuning the keypoint matchers with a small number of labels for improved precision and consistency.
%     \item investigating generalization capabilities in more detail. how does it scale with number of demos, number of objects, number of scene variations? Doing this properly most likely requires simulation to ensure the exact same demonstrations can be used with varying scene conditions etc.. cf. \cite{lin2025imitation_scaling_laws}
%     \item how to use wrist cam? would help RGB, can it help keypoints? can project onto wrist image (if visible) from same config, but does matching/tracking work as well? probably requires wide FOV camera.. 

%     \item how to combine keypoints with other representations (masks, depth, etc.) info for policy?
%     \item other ways to lift to 3d, including triangulation? can extend scope to objects that are transparent or reflective.
% \end{itemize}
\section{Conclusion}

We present an in-depth evaluation of keypoint imitation learning (KIL). KIL outperforms RGB policies under scene variations but performs similarly within distribution and for unseen objects. Additionally, KIL performs on par with S$^2$-Diffusion, raising questions about the added value of keypoints compared to other representations. For keypoint extraction, our results show that \textit{image matching }performs similarly to the more complex  \textit{instance matching} and \textit{tracking} methods. We further showed that a dedicated transformer encoder improves performance, and also extended KIL to handle tasks with multiple instances of the same object category. Finally, we illustrated how KIL inherits poor robustness to large orientation changes as a limitation from the foundation models used for keypoint matching, in addition to their limited precision and ability to handle occlusions. 
Our findings clarify several design choices for KIL and confirm its potential for improved generalization, but more research is needed to overcome its limitations and understand its added value compared to alternative representations.

\section*{ACKNOWLEDGMENTS}
\ifanonymous
Funding anonymized for double-blind review.
\else
T.L. and M.M. were funded by the European project euROBIN (grant No. 101070596); T.L. was also funded by FWO (grant No. 1S56024N).
\fi

We used LLMs (GPT5, Claude Sonnet/Opus 4.5) to assist with coding, paper formatting and proofreading.

% \textbf{generative AI}
% Thomas 
% \begin{itemize}
%     \item Claude Opus 4.5 and Sonnet 4.5/4.6 for Coding
%     \item Gemini 3(.1) pro and Claude Sonnet/Opus 4.5 for creating latex tables, brainstorming and help with proofreading.
% \end{itemize}

% Marco - ChatGPT (5.2) for coding and for helping with syntax and rephrasing in writing.



\bibliographystyle{IEEEtran}
\bibliography{references}

\ifappendixincluded

\clearpage
\onecolumn % two-col kept fighting me for whitespaces between tables etc, as it really wants to fill entire pages..

\section*{APPENDIX}
\FloatBarrier % no tables above header

% Redefine subsection/subsubsection numbering for appendix:
% \subsection -> Appendix A, B, ...
% \subsubsection -> Appendix A.1, A.2, ...
\setcounter{subsection}{0}
\renewcommand{\thesubsection}{\Alph{subsection}}
\renewcommand{\thesubsubsection}{\Alph{subsection}.\arabic{subsubsection}}


% \subsection{Contribution Statements}
% \textbf{Thomas Lips}: Conceptualization, Methodology, Software, Experiments, Writing -- Original Draft, Writing -- Review \& Editing.
% \textbf{Marco Moletta}: Conceptualization, Methodology, Software, Experiments, Writing -- Original Draft, Writing -- Review \& Editing.
% \textbf{Michael C. Welle}: Conceptualization, Methodology, Writing -- Review \& Editing, Supervision.
% \textbf{Danica Kragic}: Writing -- Review , Supervision, Funding Acquisition.
% \textbf{Francis Wyffels}: Writing -- Review, Supervision, Funding Acquisition.

\subsection{Additional Experiments}
\label{sec:appendix-additional-experiments}
In this appendix, we provide two additional experiments. In Appendix~\ref{sec:appendix-additional-experiments-augmentations}, we investigate the impact of the data augmentations described in Section~\ref{sec:method-policy-learning} on KIL performance. In Appendix~\ref{sec:appendix-additional-experiments-token-pooling}, we compare the performance of different token pooling strategies for KIL, for tasks with multiple object instances, as discussed in Section~\ref{sec:experiments-multi-objects}.
\subsubsection{Impact of Augmentations}
\label{sec:appendix-additional-experiments-augmentations}
We study the impact of data augmentations for KIL by selectively enabling the two main augmentations described in Section~\ref{sec:method-policy-learning}: adding \textit{Keypoint Noise} and applying a random \textit{Spatial Transform (ST)}. We measure performance on \textit{Pick Pen} and \textit{Place Mouse} using \textit{image matching} as keypoint extraction methods. 


Table~\ref{tab:augmentation-comparison} reports results. We find that not using any augmentation ($42\%$) hurts performance, but do not observe meaningful differences between using \textit{KP Noise} ($80\%$), \textit{ST} ($70\%$) or both ($77\%$). For all experiments in this paper, we have used both augmentations. 

At some point, we also tested keypoint dropout augmentation, as in \cite{wang2025skil}, but found it to degrade performance significantly on some tasks, hence we did not use it in the experiments presented in this paper.

\begin{table*}
    \renewcommand{\arraystretch}{1.2}
    \begin{minipage}[t]{0.38\linewidth}
        \centering
        \setlength{\tabcolsep}{2pt}
        \caption{Impact of Augmentations on KIL. Evaluated on the \textit{Pick Pen} and \textit{Place Mouse} tasks.}
        \label{tab:augmentation-comparison}
        \begin{tabularx}{\linewidth}{l Y c}
            \toprule
            \textbf{Augmentations} & \textbf{Avg.} & \textbf{CI (\%)} \\
            \midrule
            None  & 42\% & {\scriptsize [29, 55]} \\
            Noise & 80\% & {\scriptsize [68, 89]} \\
            ST    & 70\% & {\scriptsize [57, 81]} \\
            \textbf{Noise + ST}   & 77\% & {\scriptsize [64, 87]} \\
            \bottomrule
        \end{tabularx}
    \end{minipage}%
    \hfill
    \begin{minipage}[t]{0.55\linewidth}
        \centering
        \small
        \setlength{\tabcolsep}{8pt}
        \caption{Comparison of different token aggregation strategies for the keypoint encoder. Evaluated on the 10 in-distribution configurations for each task.}
        \label{table:token-pooling-comparison}
        \begin{tabularx}{\linewidth}{l Y Y Y}
            \toprule
            \textbf{Method} & \textbf{2 shoes} & \textbf{2 pens} & \textbf{pour cup} \\
            \midrule
            ours & 2 & 5.5 & 7 \\
            ours (d=256) & 2.5 & 3.5 & 8 \\
            \textbf{ours (pointBERT)} & 8 & 7.5 & 9 \\
            none & 1.5 & 3 & 6 \\
            \bottomrule
        \end{tabularx}
    \end{minipage}
\end{table*}
\subsubsection{Token Pooling Strategy Comparison}
\label{sec:appendix-additional-experiments-token-pooling}
In this experiment we elaborate on the use of a different token aggregation strategy in Section~\ref{sec:experiments-multi-objects}. As stated before, we found that using pointBERT~\cite{yu2022point-bert} token aggregation (CLS token + maxpool) worked better than using mean-pooling, which we adopted from~\cite{wang2025skil} for these tasks with multiple object instances.

We hypothesize that as the number of (independent) keypoints increases, mean-pooling dilutes spatially relevant features by a factor of 
\(\frac{1}{N}\), reducing the accuracy of the representation under mean-pooling or forcing the network to learn a more complex mapping to compensate for the dilution.

In Table~\ref{table:token-pooling-comparison}, we compare mean pooling against pointBERT pooling. We also compare against widening the transformer encoder, thus providing more capacity to the tokens, or not using an encoder at all. We find that pointBERT  outperforms all other options for the two-object tasks, confirming the importance of using an appropriate token aggregation strategy.

We also applied pointBERT pooling to the single-object tasks, as illustrated by the comparison on the \textit{pour} task in Table~\ref{table:token-pooling-comparison}. We find that it performs similarly to mean-pooling.
We expect that the performance for the single-object tasks would in general improve slightly with pointBERT pooling. However, we expect it to boost performance for all keypoint-based methods equally, hence not affecting the conclusions of the experiments in this work. We recommend future works to avoid mean-pooling or to at least compare different aggregation strategies.


\subsection{Implementation Details and Hyperparameters}
\label{sec:appendix-implementation-details-and-hyperparameters}
In this section, we provide additional details on the implementation and hyperparameters used for the policies.


Table~\ref{tab:hparams} summarises the hyperparameters used for training, except for the augmentation hyperparameters, which are reported in Table~\ref{tab:augmentations}. 

Table~\ref{tab:inference-latency} reports the approximate inference latency for each method on an RTX 4090 with torch v2.9.1. We use bf16 precision for foundation model inference and fp32 for policy inference. We do not use torch.compile. The table reports both per-observation latency (preprocessing of each observation at 10Hz, if needed) and the per-inference-step latency (the full cost paid each time the policy is queried for a chunk of actions).

Finally, in Listing~\ref{lst:encoder}, we provide a minimal implementation of the encoder architecture used for the keypoint-based policies, as described in Section~\ref{sec:method-policy-learning}.

\begin{table}[h]
    \centering
    \renewcommand{\arraystretch}{1.3}
    \setlength{\tabcolsep}{8pt}
    \caption{Hyperparameters for KIL and baseline policies. All methods use a Diffusion Policy~\cite{chi2025diffusion_policy} action head.}
    \label{tab:hparams}
    \begin{tabularx}{\linewidth}{l Y Y}
        \toprule

        \textbf{Hyperparameter} & \textbf{RGB / S$^2$} & \textbf{KIL} \\
        \midrule
        Observation encoder & ResNet18 & Transformer \\
        Input resolution    & $320\times240$ & $640\times480$ \\
        Action Head Dims                          & [256,512,1024] & [128,256,512] \\
        Action Head kernel size   & \multicolumn{2}{c}{5} \\
        Noise scheduler   & \multicolumn{2}{c}{DDIM} \\
        Denoising steps   & \multicolumn{2}{c}{100 train /14 inference} \\
        Action prediction horizon ($T_p$)   & \multicolumn{2}{c}{16} \\
        Action horizon ($T_a$)                       & \multicolumn{2}{c}{8} \\
        Observation horizon ($T_o$) & \multicolumn{2}{c}{2} \\
        Peak LR       & \multicolumn{2}{c}{$10^{-4}$} \\
        LR schedule                            & \multicolumn{2}{c}{cosine} \\
        Optimizer & Adam & AdamW \\
        Weight decay & $10^{-6}$ & $10^{-5}$ \\
        Batch size          & 64 & 128 \\
        Training steps      & 100K & 250K \\
        Augmentations (Table~\ref{tab:augmentations})       & RandomCrop, Gaussian Noise, ColorJitter & KP Noise, ST, EEF Noise \\
        \bottomrule
    \end{tabularx}
\end{table}

\begin{table*}[h]
    \renewcommand{\arraystretch}{1.3}
    \begin{minipage}[t]{0.48\linewidth}
        \centering
        \setlength{\tabcolsep}{8pt}
        \caption{Augmentation hyperparameters. ST = Spatial Transform; KP = Keypoint; EEF = End-Effector.}
        \label{tab:augmentations}
        \begin{tabularx}{\linewidth}{l l X}
            \toprule
            \textbf{Method} & \textbf{Augmentation} & \textbf{Parameters} \\
            \midrule
            \multirow{3}{*}{RGB}
                & RandomCrop & $288\times216$ \\
                & Gaussian Noise & $\sigma = 12.5$ mm \\
                & ColorJitter & range$=(0.4, 0.4, 0.4, 0.1)$ \\
            \midrule
            \multirow{2}{*}{S$^2$}
                & RandomCrop & $288\times216$ \\
                & Gaussian Noise & $\sigma = 12.5$ px \\
            \midrule
            \multirow{3}{*}{KIL}
                & KP Noise & $\sigma = 1$ mm \\
                & ST & $\Delta t = 0.3$ m, $\Delta\theta_z = 0.6$ rad \\
                & EEF Noise & $\sigma_{\text{pos}} = 2$ mm, $\sigma_{\text{ori}} = 1\deg $, $\sigma_{\text{grip}} = 0.02$ \\
            \bottomrule
        \end{tabularx}
    \end{minipage}%
    \hfill
    \begin{minipage}[t]{0.49\linewidth}
        \centering
        \renewcommand{\arraystretch}{1.2}
        \setlength{\tabcolsep}{3pt}
        \caption{Inference latency (ms) of the different policies used in this work, measured on an NVIDIA RTX 4090 using bf16 precision for foundation model inference. We do not use torch.compile.}
        \label{tab:inference-latency}
        \begin{tabularx}{\linewidth}{l c Y}
            \toprule
            \textbf{Method} & \textbf{Per observation} & \textbf{Per inference step} \\
            \midrule
            RGB              & /  & $N_o \times$ ResNet18 + AH : 35 ms \\
            S2               &/ & $N_o \times$ (DAv2 + gSAM) + AH: 540 ms \\
            KIL (IM) & /    & $N_o \times$ RADIOv2.5-B + AH: 80 ms \\
            KIL (IN) & /  & $N_o \times$ (DIFT\textsubscript{2} + SAM3) + AH: 570 ms\\
            KIL (T)    & CoTracker3 : 59 ms           & AH : 29ms \\
            \bottomrule
        \end{tabularx}
    \end{minipage}
\end{table*}




\begin{figure}[h]
\noindent\begin{minipage}{\linewidth}
\begin{lstlisting}[language=Python, caption=Minimal implementation of the keypoint encoder, label=lst:encoder,
  basicstyle=\ttfamily\scriptsize, breaklines=false, columns=fullflexible]
  import torch, torch.nn as nn, numpy as np

  def sincos_embed(d, n):
      i = np.arange(d // 2) / (d / 2.0)
      e = np.outer(np.arange(n), 1.0 / 10000**i)
      return torch.from_numpy(np.concatenate([np.sin(e), np.cos(e)], axis=1)).float()
  
  class KeypointEncoder(nn.Module):
      def __init__(self, seq_len=9, kp_dim=3, ee_dim=7,d=128, nhead=8, num_layers=4, ff=512):
          super().__init__()
          self.seq_len = seq_len
          self.kp_dim  = kp_dim
          self.ee_dim  = ee_dim
          self.d       = d
  
          self.kp_emb = nn.Linear(kp_dim, d)
          self.ee_mlp = nn.Sequential(
              nn.Linear(ee_dim, d*2), nn.ReLU(),
              nn.Linear(d*2,   d*2), nn.ReLU(),
              nn.Linear(d*2,   d),   nn.Tanh(),
          )
          layer = nn.TransformerEncoderLayer(d, nhead, ff, dropout=0.0,activation='gelu', batch_first=True,norm_first=True)
          self.encoder = nn.TransformerEncoder(layer, num_layers, norm=None)
          self.register_buffer('pos', sincos_embed(d, seq_len + 1).unsqueeze(0))
  
      def forward(self, kp, ee):
          B, T = kp.shape[:2]
          kp_tokens = self.kp_emb(kp.reshape(B*T, self.seq_len, self.kp_dim)) + self.pos[:, :self.seq_len]
          ee_token  = self.ee_mlp(ee.reshape(B*T, self.ee_dim)).unsqueeze(1) + self.pos[:, self.seq_len:]
          tokens = torch.cat([kp_tokens, ee_token], dim=1)
          tokens = self.encoder(tokens)
          return tokens.mean(dim=1).reshape(B, T, -1)
\end{lstlisting}
\end{minipage}
\end{figure}


% \subsection{Art Is in the Details}
% Some details/findings/lessons learned:
% Continuous vs discrete gripper state. discrete ( asymmetric) would have boosted performance because now robot is sometimes in unknown state-space when sensor says stop but gripper is still kind of open.


% \subsection{Hardware Is Challenging}
% \thomas{should probably not be in the paper, but want it in my dissertation so this serves as a reminder.}

% Camera frame corruption issue.

% Gripper bug.

% Both hard to detect. in both cases I first thought it was artifact of the policy. 

\subsection{Initial Scene Configurations of Demonstrations}
\label{sec:appendix-demo-configurations}
Figure~\ref{fig:demo-distributions} shows an overlay of the initial frames of the demonstrations for each task, illustrating the distribution of initial scene configurations used during data collection. As described in Section~\ref{sec:experiments-tasks}, we vary the position of objects significantly across the workspace, but limit orientation variation to $\pm 30^\circ$ around a canonical orientation. Demonstrations are collected using two different objects per task, without scene variations.

\begin{figure}[H]
    \centering
    \subfloat[Place Shoe]{\includegraphics[width=0.32\linewidth]{imgs/demo_distribution/1shoe-v4-distribution.png}}%
    \hfill
    \subfloat[Place Mug]{\includegraphics[width=0.32\linewidth]{imgs/demo_distribution/1mug-v3-distribution.png}}%
    \hfill
    \subfloat[Pour Bottle]{\includegraphics[width=0.32\linewidth]{imgs/demo_distribution/pour-bottle-v1-distribution.png}}\\
    \subfloat[Pick Pen]{\includegraphics[width=0.32\linewidth]{imgs/demo_distribution/1pen-distribution.png}}%
    \hfill
    \subfloat[Place Mouse]{\includegraphics[width=0.32\linewidth]{imgs/demo_distribution/1mouse-distribution.png}}%
    \hfill
    \subfloat{\phantom{
\includegraphics[width=0.32\linewidth]{imgs/demo_distribution/black.png}}}%
    \vspace{4pt}
    \caption{Scene configuration distribution of the demonstrations for each task.}
    \label{fig:demo-distributions}
\end{figure}


\subsection{Keypoint Representation Reference images}
\label{sec:appendix-keypoint-references}
Figure~\ref{fig:kp-configs} shows the keypoint reference images used for each task, as described in Section~\ref{sec:method-representation-specification}. Each image shows the reference frame with the annotated keypoints and text prompt used for instance detection (cf. Section~\ref{sec:method-representation-extraction}).

\begin{figure}[H]
    \centering
    \includegraphics[width=0.49\linewidth]{imgs/kp_configs/1shoe-tray-v2-kp.png}\hfill
    \includegraphics[width=0.49\linewidth]{imgs/kp_configs/1mug-tray-v5-kp.png}\\[8pt]
    \includegraphics[width=0.49\linewidth]{imgs/kp_configs/pour-bottle-v1-kp.png}\hfill
    \includegraphics[width=0.49\linewidth]{imgs/kp_configs/1mouse-both-7kp.png}\\[8pt]
    \includegraphics[width=0.49\linewidth]{imgs/kp_configs/1pen-both-7kp.png}
    \caption{Keypoint reference images for each task, inluding text prompts for each object.}
    \label{fig:kp-configs}
\end{figure}

\subsection{Initial Scene Configurations for Evaluations}
\label{sec:appendix-eval-configurations}

Figure~\ref{fig:eval-configurations} shows the initial scene configurations used for evaluation. For each task, we measure in-distrubtion performance, generalization to unseen objects and generalization to scene variations. We use the same 10 initial scene configurations for all policies to ensure fairness.

\begin{figure*}[h]
    \centering
    \begin{minipage}[c]{0.03\linewidth}\centering\rotatebox{90}{Place Shoe}\end{minipage}%
    \begin{minipage}[c]{0.97\linewidth}\includegraphics[width=\linewidth]{imgs/eval_configs/1shoe-v4-grid.jpg}\end{minipage}\\[6pt]
    \begin{minipage}[c]{0.03\linewidth}\centering\rotatebox{90}{Place Mug}\end{minipage}%
    \begin{minipage}[c]{0.97\linewidth}\includegraphics[width=\linewidth]{imgs/eval_configs/1mug-v3-grid.jpg}\end{minipage}\\[6pt]
    \begin{minipage}[c]{0.03\linewidth}\centering\rotatebox{90}{Pour Bottle}\end{minipage}%
    \begin{minipage}[c]{0.97\linewidth}\includegraphics[width=\linewidth]{imgs/eval_configs/pour-bottle-v1-grid.jpg}\end{minipage}\\[6pt]
    \begin{minipage}[c]{0.03\linewidth}\centering\rotatebox{90}{Pick Pen}\end{minipage}%
    \begin{minipage}[c]{0.97\linewidth}\includegraphics[width=\linewidth]{imgs/eval_configs/1pen-grid.jpg}\end{minipage}\\[6pt]
    \begin{minipage}[c]{0.03\linewidth}\centering\rotatebox{90}{Place Mouse}\end{minipage}%
    \begin{minipage}[c]{0.97\linewidth}\includegraphics[width=\linewidth]{imgs/eval_configs/1mouse-grid.jpg}\end{minipage}\\[6pt]
    \caption{Initial scene configurations used for evaluation. For each task, we use 10 in-distribution configurations (first row), 10 unseen object configurations (second row) and 10 scene variation configurations (third row).}
    \label{fig:eval-configurations}
\end{figure*}


\subsection{Unaggregated Results}
\label{sec:appendix-detailed-results}
In this appendix, we provide the unaggregated results for tables in the main text and in Appendix~\ref{sec:appendix-additional-experiments}, broken down by task and evaluation setting (in-distribution, unseen objects, scene variations). We only include tables for which the unaggregated results were not already shown.


\begin{table*}[t]
    \renewcommand{\arraystretch}{1.2}
    \begin{minipage}[t]{0.63\linewidth}
        \centering
        \setlength{\tabcolsep}{4pt}
        \caption{Unaggregated results for Table~\ref{tab:model-comparison}: Comparison of keypoint feature models for the different keypoint extraction methods. The models used in the other experiments in this paper are marked in bold.}
        \label{tab:model-comparison-full}
        \begin{tabularx}{\linewidth}{l *{3}{Y} *{3}{Y} {c}}
            \toprule
            \multirow{2}{*}{\textbf{Method}} & \multicolumn{3}{c}{\textbf{place shoe}} & \multicolumn{3}{c}{\textbf{place mug}} & \multirow{2}{*}{\textbf{Total (/60)}} \\
            \cmidrule(lr){2-4} \cmidrule(lr){5-7}
            & in & obj & scene & in & obj & scene & \\
            \midrule
            RGB~\cite{chi2025diffusion_policy} & 8 & 6 & 3 & 10 & 7 & 0 & 34 \\
            \midrule
            \textbf{RADIOv2.5-B}  & 10 & 6 & 9 & 10 & 8 & 8 & 51 \\
            DIFT$_2$  & 7 & 4 & 2 & 10 & 8& 6 &37  \\
            DINOv3-B & 10 & 4 & 6 & 10 & 5 & 9 & 44  \\
            \midrule
            Instance + RADIOv2.5-B &7 & 6& 6&9 &8 &10 & 46 \\
            \textbf{Instance + DIFT$_2$} & 10 & 6 & 5 & 10 & 7 & 10 & 48 \\
            \midrule
            Tracking + RADIOv2.5-B & 7&6 &7 &9 &6 & 9 & 44 \\
            Tracking + DIFT$_2$ & 9 & 8 & 6 & 10 & 7 & 9 & 49 \\
            \textbf{Tracking + DIFT$_8$} &9 & 6&6 &9 & 9 & 9 & 48 \\
            \bottomrule
        \end{tabularx}
    \end{minipage}%
    \hfill
    \begin{minipage}[t]{0.34\linewidth}
        \centering
        \setlength{\tabcolsep}{5pt}
        \caption{Unaggregated results for Table~\ref{tab:augmentation-comparison}: Impact of Augmentations on KIL.}
        \label{tab:augmentation-comparison-full}
        \begin{tabularx}{\linewidth}{l *{3}{Y} *{3}{Y} {c}}
            \toprule
            \multirow{2}{*}{\textbf{Method}} & \multicolumn{3}{c}{\textbf{pick pen}} & \multicolumn{3}{c}{\textbf{place mouse}} & \multirow{2}{*}{\textbf{Total (/60)}} \\
            \cmidrule(lr){2-4} \cmidrule(lr){5-7}
            & in & obj & scene & in & obj & scene & \\
            \midrule
            None & 7 & 10  & 3 & 1 & 2 & 2 &  25 \\
            Noise  & 6 & 9  & 7 & 10 & 8 & 8 & 48  \\
            ST & 5 & 6  & 4 & 10 & 8 & 9 & 42  \\
            Noise + ST & 7 & 8  & 4 & 10 & 8 & 9 &  46 \\
            \bottomrule
        \end{tabularx}
        \vspace{1em}
        \setlength{\tabcolsep}{2pt}
        \caption{Unaggregated results for Table~\ref{tab:encoder-comparison}: Encoder Comparison.}
        \label{table:encoder-comparison-full}
        \begin{tabularx}{\linewidth}{l *{3}{Y} *{3}{Y} c}
            \toprule
            \multirow{2}{*}{\textbf{Method}} & \multicolumn{3}{c}{\textbf{place shoe}} & \multicolumn{3}{c}{\textbf{place mug}} & \multirow{2}{*}{\textbf{Total (/60)}} \\
            \cmidrule(lr){2-4} \cmidrule(lr){5-7}
            & in & obj & scene & in & obj & scene & \\
            \midrule
            \textbf{ours} & 10 & 6 & 9 & 10 & 8 & 8 & \textbf{51} \\
            None   & 10 & 9 & 9 & 8 & 7 & 8 & \textbf{51} \\
            \hspace{.3em} + sim   & 2 & 0 & 2 & 0 & 0 & 0 & 4 \\
            \hspace{1.2em} - aug & 7 & 6 & 7 & 9 & 5 & 5 & 39 \\
            \bottomrule
        \end{tabularx}
    \end{minipage}
\end{table*}

\begin{table*}[t]
    \renewcommand{\arraystretch}{1.2}
    \begin{minipage}[t]{0.54\linewidth}
        \centering
        \small
        \setlength{\tabcolsep}{5pt}
        \caption{Unaggregated results for Table~\ref{table:2-object-experiments}: Task completion for tasks in which multiple instances of the same object category are handled.}
        \label{table:2-object-experiments-full}
        \begin{tabularx}{\linewidth}{l *{6}{Y} c}
            \toprule
            \multirow{2}{*}{\textbf{Method}} &
            \multicolumn{3}{c}{\textbf{place 2 shoes}} &
            \multicolumn{3}{c}{\textbf{pick 2 pens}} &
            \multirow{2}{*}{\textbf{Total (/60)}} \\
            \cmidrule(lr){2-4} \cmidrule(lr){5-7}
            & in & obj & scene & in & obj & scene & \\
            \midrule
            RGB~\cite{chi2025diffusion_policy}      & 6.5 & 4   & 2.5 & 8.5 & 4.5 & 2.5 & 28.5 \\
            S$^2$~\cite{yang2025S2}       & 4   & 5.5 & 5   & 8   & 7   & 5.5 & 35 \\
            KIL (IM)       & 8.5 & 6   & 7   & 7.5 & 4.5 & 4.5 & 38 \\
            KIL (IN) & 8 & 7.5 & 6 & 7.5 & 5.5 & 2  & 36.5 \\
            KIL (T) & 7.5 & 4.5 & 3.5 & 6 & 5.5 & 4 & 31 \\
            \bottomrule
        \end{tabularx}
    \end{minipage}%
    \hfill
    \begin{minipage}[t]{0.43\linewidth}
        \centering
        \setlength{\tabcolsep}{5pt}
        \caption{Unaggregated results for Table~\ref{tab:more-rotations-experiment}: Illustration of performance drop when increasing the range of initial orientations for task objects.}
        \label{tab:more-rotations-experiment-full}
        \begin{tabularx}{\linewidth}{l*{3}{Y}*{3}{Y}c}
            \toprule
            \multirow{2}{*}{\textbf{Method}} & \multicolumn{3}{c}{\textbf{pick pen ($\pm 180^\circ$)}} & \multicolumn{3}{c}{\textbf{place mouse ($\pm 180^\circ$)}} & \multirow{2}{*}{\makecell{\textbf{Total}\\\textbf{(/60)}}} \\
            \cmidrule(lr){2-4} \cmidrule(lr){5-7}
            & in & obj & scene & in & obj & scene & \\
            \midrule
            RGB & 6 & 7  & 0 & 8 & 7 & 0 &  28 \\
            S$^2$  & 7 & 6  & 5 & 9 & 10 & 5 & \textbf{42} \\
            KIL (IM) & 6 & 4  & 2 & 5 & 3 & 1 & 21 \\
            KIL (T) & 6 & 5  & 4 & 5 & 3 & 6 &  29 \\
            \bottomrule
        \end{tabularx}
    \end{minipage}
\end{table*}


\fi % ifappendixincluded

\end{document}

% \subsection{Unaggregated Results}
% \label{sec:appendix-detailed-results}
% In this appendix, we provide the unaggregated results for tables in the main text and in Appendix~\ref{sec:appendix-additional-experiments}, broken down by task and evaluation setting (in-distribution, unseen objects, scene variations). We only include tables for which the unaggregated results were not already shown.


% \begin{table*}[t]
%     \centering
%     \renewcommand{\arraystretch}{1.2}
%     \setlength{\tabcolsep}{4pt}
%     \caption{Unaggregated results for Table~\ref{tab:model-comparison}: Comparison of keypoint feature models for the different keypoint extraction methods. The models used in the other experiments in this paper are marked in bold.}
%     \label{tab:model-comparison-full}
%     \begin{tabularx}{\linewidth}{l *{3}{Y} *{3}{Y} {c}}
%         \toprule
%         \multirow{2}{*}{\textbf{Method}} & \multicolumn{3}{c}{\textbf{place shoe}} & \multicolumn{3}{c}{\textbf{place mug}} & \multirow{2}{*}{\textbf{Total (/60)}} \\
%         \cmidrule(lr){2-4} \cmidrule(lr){5-7}
%         & in & obj & scene & in & obj & scene & \\
%         \midrule
%         RGB~\cite{chi2025diffusion_policy} & 8 & 6 & 3 & 10 & 7 & 0 & 34 \\
%         \midrule
%         \textbf{RADIOv2.5-B}  & 10 & 6 & 9 & 10 & 8 & 8 & 51 \\
%         DIFT$_2$  & 7 & 4 & 2 & 10 & 8& 6 &37  \\
%         DINOv3-B & 10 & 4 & 6 & 10 & 5 & 9 & 44  \\
%         \midrule
%         Instance + RADIOv2.5-B &7 & 6& 6&9 &8 &10 & 46 \\
%         \textbf{Instance + DIFT$_2$} & 10 & 6 & 5 & 10 & 7 & 10 & 48 \\
%         \midrule
%         Tracking + RADIOv2.5-B & 7&6 &7 &9 &6 & 9 & 44 \\
%         Tracking + DIFT$_2$ & 9 & 8 & 6 & 10 & 7 & 9 & 49 \\
%         \textbf{Tracking + DIFT$_8$} &9 & 6&6 &9 & 9 & 9 & 48 \\
%         \bottomrule
%     \end{tabularx}
% \end{table*}

% \begin{table*}[t]
%     \centering
%     \renewcommand{\arraystretch}{1.2}
%     \setlength{\tabcolsep}{5pt}
%     \caption{Unaggregated results for Table~\ref{tab:augmentation-comparison}: Impact of Augmentations on KIL.}
%     \label{tab:augmentation-comparison-full}
%     \begin{tabularx}{\linewidth}{l *{3}{Y} *{3}{Y} {c}}
%         \toprule
%         \multirow{2}{*}{\textbf{Method}} & \multicolumn{3}{c}{\textbf{pick pen}} & \multicolumn{3}{c}{\textbf{place mouse}} & \multirow{2}{*}{\textbf{Total (/60)}} \\
%         \cmidrule(lr){2-4} \cmidrule(lr){5-7}
%         & in & obj & scene & in & obj & scene & \\
%         \midrule
%         None & 7 & 10  & 3 & 1 & 2 & 2 &  25 \\
%         Noise  & 6 & 9  & 7 & 10 & 8 & 8 & 48  \\
%         ST & 5 & 6  & 4 & 10 & 8 & 9 & 42  \\
%         Noise + ST & 7 & 8  & 4 & 10 & 8 & 9 &  46 \\
%         \bottomrule
%     \end{tabularx}
% \end{table*}

% \begin{table*}[t]
%     \centering
%     \renewcommand{\arraystretch}{1.2}
%     \setlength{\tabcolsep}{2pt}
%     \caption{Unaggregated results for Table~\ref{tab:encoder-comparison}: Encoder Comparison.}
%     \label{table:encoder-comparison-full}
%     \begin{tabularx}{\linewidth}{l *{3}{Y} *{3}{Y} c}
%         \toprule
%         \multirow{2}{*}{\textbf{Method}} & \multicolumn{3}{c}{\textbf{place shoe}} & \multicolumn{3}{c}{\textbf{place mug}} & \multirow{2}{*}{\textbf{Total (/60)}} \\
%         \cmidrule(lr){2-4} \cmidrule(lr){5-7}
%         & in & obj & scene & in & obj & scene & \\
%         \midrule
%         \textbf{ours} & 10 & 6 & 9 & 10 & 8 & 8 & \textbf{51} \\
%         None   & 10 & 9 & 9 & 8 & 7 & 8 & \textbf{51} \\
%         \hspace{.3em} + sim   & 2 & 0 & 2 & 0 & 0 & 0 & 4 \\
%         \hspace{1.2em} - aug & 7 & 6 & 7 & 9 & 5 & 5 & 39 \\
%         \bottomrule
%     \end{tabularx}
% \end{table*}

% \begin{table*}[t]
%     \centering
%     \renewcommand{\arraystretch}{1.2}
%     \small
%     \setlength{\tabcolsep}{5pt}
%     \caption{Unaggregated results for Table~\ref{table:2-object-experiments}: Task completion for tasks in which multiple instances of the same object category are handled.}
%     \label{table:2-object-experiments-full}
%     \begin{tabularx}{\linewidth}{l *{6}{Y} c}
%         \toprule
%         \multirow{2}{*}{\textbf{Method}} &
%         \multicolumn{3}{c}{\textbf{place 2 shoes}} &
%         \multicolumn{3}{c}{\textbf{pick 2 pens}} &
%         \multirow{2}{*}{\textbf{Total (/60)}} \\

%         \cmidrule(lr){2-4} \cmidrule(lr){5-7}

%         & in & obj & scene & in & obj & scene & \\
%         \midrule
%         RGB~\cite{chi2025diffusion_policy}      & 6.5 & 4   & 2.5 & 8.5 & 4.5 & 2.5 & 28.5 \\
%         S$^2$~\cite{yang2025S2}       & 4   & 5.5 & 5   & 8   & 7   & 5.5 & 35 \\
%         KIL (IM)       & 8.5 & 6   & 7   & 7.5 & 4.5 & 4.5 & 38 \\
%         KIL (IN) & 8 & 7.5 & 6 & 7.5 & 5.5 & 2  & 36.5 \\
%         KIL (T) & 7.5 & 4.5 & 3.5 & 6 & 5.5 & 4 & 31 \\
%         \bottomrule
%     \end{tabularx}
% \end{table*}

% \begin{table*}[t]
%     \centering
%     \renewcommand{\arraystretch}{1.2}
%     \setlength{\tabcolsep}{5pt}
%     \caption{Unaggregated results for Table~\ref{tab:more-rotations-experiment}: Illustration of performance drop when increasing the range of initial orientations for task objects.}
%     \label{tab:more-rotations-experiment-full}
%     \begin{tabularx}{\linewidth}{l*{3}{Y}*{3}{Y}c}
%         \toprule
%         \multirow{2}{*}{\textbf{Method}} & \multicolumn{3}{c}{\textbf{pick pen ($\pm 180^\circ$)}} & \multicolumn{3}{c}{\textbf{place mouse ($\pm 180^\circ$)}} & \multirow{2}{*}{\makecell{\textbf{Total}\\\textbf{(/60)}}} \\
%         \cmidrule(lr){2-4} \cmidrule(lr){5-7}
%         & in & obj & scene & in & obj & scene & \\
%         \midrule
%         RGB & 6 & 7  & 0 & 8 & 7 & 0 &  28 \\
%         S$^2$  & 7 & 6  & 5 & 9 & 10 & 5 & \textbf{42} \\
%         KIL (IM) & 6 & 4  & 2 & 5 & 3 & 1 & 21 \\
%         KIL (T) & 6 & 5  & 4 & 5 & 3 & 6 &  29 \\
%         \bottomrule
%     \end{tabularx}
% \end{table*}
% \vspace*{\fill}


% \fi % ifappendixincluded

% \end{document}